% \documentclass[10pt]{article}
% \usepackage[french]{babel}
% \usepackage[utf8]{inputenc}
% \usepackage[T1]{fontenc}
% \usepackage{amsmath}
% \usepackage{amsfonts}
% \usepackage{amssymb}
% \usepackage[version=4]{mhchem}
% \usepackage[left=2cm,right=2cm,top=2cm,bottom=2cm]{geometry}
% \usepackage{stmaryrd}
\input{TD.sty}
\newcommand{\pcsi}[1]{
  \begin{liste}
    \item \udl{\textsc{PCSI}}~:~#1
  \end{liste}
}
\newcommand{\etudiant}[2][*]{
  \begin{itemize}
\ifthenelse{\equal{\detokenize{#1}}{\detokenize{*}}}
{\item \color{red}\fbox{\uppercase{\textbf{#2}}}~:}
{\item \fbox{\uppercase{\textbf{#2}}}~:}
\end{itemize}
}

\newcommand{\pc}[1]{
  \begin{liste}
\item \fbox{\uppercase{\textbf{pc}}}~:~#1
\end{liste}
}


\begin{document}
\begin{center}
  \Large
  \textsc{TIPE de physique}
\end{center}

\etudiant{MOURET - BAMAS}
\begin{liste}
  \item 05/09/2024 : Sujet : \udl{Barrage hydraulique} ou \udl{Turbine} ? Idée : produire un modèle réduit d’un barrage. Étudier le rendement, la résistance.
  \item 12/09/2024 : 
  \item 19/09/2024 : CHIMIE : \udl{Pile}.
  \item 26/09/2024 : test de la pile : débit dans un résistance.
\end{liste}

\etudiant{MORLAN - CASTET}
\begin{liste}
  \item 05/09/2024 : Sujet : \udl{Train d’aterrissage} : conversion en frottements ? \udl{Catalyseur de voiture} : plutôt de la chimie ?
  \item 12/09/2024 : 
  \item 19/09/2024 : Sujet : \udl{Tomographie laser}.
  \item 26/09/2024 : \udl{Modulation} de signal.
\end{liste}

\etudiant{BRAUX - SANJAUME}
\begin{liste}
  \item 05/09/2024 : Sujet : pas de sujet !
  \item 12/09/2024 : 
  \item 19/09/2024 : 
  \item 26/09/2024 : ?
\end{liste}

\etudiant{SCELLES - BICHEREL} 
\begin{liste}
  \item 05/09/2024 : 
  Sujet : \udl{métro de Paris} : étude rail électrique, frottements. Est bien renseigné là-dessus. Expérience ?
  \item 12/09/2024 : 
  \item 19/09/2024 : 
  \item 26/09/2024 : ?
  \item  03/10/2024 : réseau électrique ? 

\end{liste}

\etudiant[]{BUISSON - }
\begin{liste}
  \item 05/09/2024 : Sujet : \udl{Tsunami} : lien avec thème de l’année ? Paramètres = forme de digue. Expérience quantitative ?
  \item 12/09/2024 : 
  \item 19/09/2024 : 
\end{liste}

\etudiant[]{BICHEREL - } 
\begin{liste}
  \item 05/09/2024 : Sujet : \udl{Tir à l’arc}. Variation de la forme des branches. Expériences déjà réalisées.
  \item 12/09/2024 : 
  \item 19/09/2024 : 
  \item 26/09/2024 :
  \item 03/10/2024 : expérience toujours pas démarrée.
  \item 17/10/2024 : sujet abandonné.
\end{liste}

\etudiant{FOURCADE - DIOT}
\begin{liste}
  \item 05/09/2024 : Sujet : \udl{force Houlomotrice}. Pas d’expérience encore. Idée : faire un modèle réduit.
  \item 12/09/2024 : 
  \item 19/09/2024 : Sujet : \udl{Isolation thermique}.
  \item 26/09/2024 : idem.
\end{liste}

\etudiant{LEMAÎTRE - BUISSON} 
\begin{liste}
  \item 05/09/2024 : Sujet : \udl{Bateau et voile}  ?
  \item 12/09/2024 : Sujet : \udl{Cellule photovoltaïque}.
  \item 19/09/2024 : Sujet : \udl{Stirling}.
  \item 26/09/2024 : sujet conservé.
  \item 17/10/2024 : autre sujet : chimie des pigments.
\end{liste}

\etudiant{LALANDE - } 
\begin{liste}
  \item 05/09/2024 : Sujet : \udl{Looping parc attraction}  ? Exp. en Lego réalisées.
  \item 12/09/2024 : 
  \item 19/09/2024 : 
  \item 26/09/2024 : ?
  \item 17/10/2024 : expérience à la maison, il faut la ramener ici !
\end{liste}

\etudiant{CHAPPAZ - LESTRADE} 
\begin{liste}
  \item 05/09/2024 : 
  \item 12/09/2024 : 
  \item 19/09/2024 : Sujet : \udl{Lévitation magnétique}.
  \item 26/09/2024 : \udl{Module Peltier} ?
  \item 17/10/2024 : Module Peltier oublié ! Biblio.
\end{liste}

\end{document}